% Part: incompleteness
% Chapter: incompleteness-provability
% Section: provability-conditions

\documentclass[../../../include/open-logic-section]{subfiles}

\begin{document}

\olfileid{inc}{inp}{prc}

\olsection{$\Th{PA}$ 的可推导性条件}

Peano算术（Peano arithmetic），或称 $\Th{PA}$，是对 $\Th{Q}$ 为所有公式附加归纳公理而扩张的理论。换言之，对于每个公式~$!A$，向 $\Th{Q}$ 中添加形如
\[
(!A(0) \land \lforall[x][(!A(x) \lif !A(x'))]) \lif \lforall[x][!A(x)]
\]
的公理。注意这实际上是一个\emph{公理模式}，也就是说，有无穷多条公理（而且事实证明 $\Th{PA}$ 并\emph{非}有限可公理化的）。但由于可以有效地判定一个符号串是否是归纳公理的一个实例，因此 $\Th{PA}$ 的公理集是可计算的。$\Th{PA}$ 是比~$\Th{Q}$ 强健得多的理论。例如，利用通常的归纳法，容易证明加法和乘法满足交换律。事实上，大多数有限性的数论与组合论证都可以在~$\Th{PA}$ 中完成。

由于 $\Th{PA}$ 是可计算公理化的，其可推导性谓词 $\Prf[\Th{PA}](x,y)$ 是可计算的，因而在~$\Th{Q}$（进而在~$\Th{PA}$）中可表示。如前所述，我们将用 $\OPrf[\Th{PA}](x,y)$ 来表示该关系的公式。令 $\OProv[\Th{PA}](y)$ 为公式 $\lexists[x][\Prf[\Th{PA}](x,y)]$，该公式直观上表示“$y$ 可由 $\Th{PA}$ 的公理推导出来”。我们之所以需要比 $\Th{Q}$ 的公理更强一点的理论，是因为我们需要确保所用的理论足够强，能够推导出关于这个可推导性谓词的一些基本事实。事实上，我们需要的是以下事实：

\begin{enumerate}
\item[P1.] 若 $\Th{PA} \Proves !A$，则 $\Th{PA} \Proves
  \OProv[\Th{PA}](\gn{!A})$。
\item[P2.] 对所有公式 $!A$ 和 $!B$，
  \[
  \Th{PA} \Proves \OProv[\Th{PA}](\gn{!A \lif !B}) \lif
  (\OProv[\Th{PA}](\gn{!A}) \lif \OProv[\Th{PA}](\gn{!B})).
  \]
\item[P3.] 对每个公式~$!A$，
  \[
  \Th{PA} \Proves \OProv[\Th{PA}](\gn{!A})
  \lif \OProv[\Th{PA}](\gn{\OProv[\Th{PA}](\gn{!A})}).
  \]
\end{enumerate}

要验证这三个性质成立，唯一的方法是仔细描述公式 $\OProv[\Th{PA}](y)$，并利用 $\Th{PA}$ 的公理来描述相关的形式推导。条件 (1) 和~(2) 是容易的；真正需要下功夫的是条件~(3)。（想想这会涉及什么样的工作……）展开所有细节会既繁琐又无趣，因此这里我们请你暂且相信 $\Th{PA}$ 具有上述列出的三个性质。对 $\OProv[\Th{PA}](y)$ 的合理选择还会满足

\begin{enumerate}
\item[P4.] 若 $\Th{PA} \Proves \OProv[\Th{PA}](\gn{!A})$，则
  $\Th{PA} \Proves !A$。
\end{enumerate}

不过我们用不上这个事实。

\begin{digress}
顺便提一下，Gödel当时也和我们现在一样偷懒。在那篇1931年的论文结尾，他概述了第二不完全性定理的证明，并承诺会在后续论文中给出细节。他后来一直没有来得及做这件事；因为所有理解该论证的人都相信它是可以完成的（他并不需要补充这些细节）。
\end{digress}

\end{document}